\documentclass[aps,prd,a4paper,10pt, twocolumn,amsmath,showpacs,superscriptaddress,nofootinbib,preprintnumbers,notitlepage]{revtex4-2}
\usepackage{dcolumn}
\usepackage{epstopdf}
\usepackage{bm}
\usepackage{braket}
\usepackage{enumitem}
\usepackage{soul}
\usepackage{pifont}
\usepackage[labelfont= bf, skip=5pt, font=small]{caption}
\usepackage{graphicx, wrapfig, subcaption}
\usepackage{array, tabularx, tabulary, booktabs, multirow}
\usepackage{textcomp, cmap, comment, epigraph}
\usepackage[dvipsnames]{xcolor}
\usepackage{orcidlink}
\usepackage[utf8]{inputenc}
\usepackage[spanish, es-tabla]{babel}
\usepackage[version=4]{mhchem}
\usepackage{hyperref}
\usepackage{geometry}
\usepackage{amsfonts, amssymb, amsthm, mathtools} 
\usepackage{amsmath}
\usepackage{float}
\usepackage{placeins}
\usepackage{siunitx}

\geometry{left=15mm, right=15mm, top=16mm, bottom=18mm}
\DeclareMathSizes{12}{16}{12}{10}
\selectlanguage{spanish}

\makeatletter
\newcommand{\fmarki}{*}
\newcommand{\fmarkii}{\ensuremath{\dagger}}
\newcommand{\fmarkiii}{\ensuremath{\ddagger}}  
\def\@fnsymbol#1{{\ifcase#1\or \fmarki\or \fmarkii\or \fmarkiii\else\@ctrerr\fi}}
\makeatother

\renewcommand{\thesection}{\Roman{section}}
\renewcommand{\abstractname}{}

\addto\captionsspanish{\renewcommand{\abstractname}{}}
\def\restrict#1{\raise-.5ex\hbox{\ensuremath|}_{#1}}
\def\sen{\mathop{\mbox{\normalfont sen}}\nolimits}
\graphicspath{ {images/} }
\hypersetup{colorlinks=true,linkcolor=MidnightBlue,citecolor=OliveGreen,urlcolor=RoyalBlue,pdftitle={Relación entre las propiedades magnéticas y geometría con el factor de avance en acero irradiado},pdfauthor={Nicolás Clotta}}

\begin{document}

\title{\LARGE Relación entre las propiedades magnéticas y geometría con~el factor de avance en acero irradiado}

\author{Nicolás Clotta}
\email{nicolas.clotta@gmail.com}
\affiliation{Departamento de Física, Facultad de Ciencias Exactas y Naturales, Universidad de Buenos Aires.}

\date[Fechado: ]{\today}
\begin{abstract}
    \noindent
\end{abstract}

\maketitle

\section{Introducción}
El acero utilizado en la construcción de recipientes de presión de reactores nucleares sufre una constante irradiación de neutrones, que junto con la elevada temperatura, eventualmente llevan al material al fallo \cite{}. Para evitar esto, resulta fundamental poder analizar el estado del mismo para intervenir a tiempo. El método usual es el test de Charpy, el cual involucra tomar muestras e impactarlas con un péndulo para así estudiar su resistencia material luego de la irradiación \cite{}. Mediante magnetometría resulta posible establecer una correlación entre el factor de avance, el flujo de neutrones que recibió la muestra por unidad de tiempo, y las propiedades magnéticas de la misma \cite{}.

\section{Métodos y materiales}

\section{Resultados y discusión}

\section{Conclusión}

\bibliographystyle{plain}
\bibliography{cuaderno6}




\end{document}
